\subsection*{Future Work}

Further research is therefore needed to detect specific behaviours from the entropy of the system more precisely. A particularly promising direction is the design and training of a machine learning model on entropy trajectories that distinguishes communication from other behavioural changes. Beyond the output distribution, a complementary direction is to exploit the internal activations of the models: lightweight probes trained on the hidden states of intermediate layers have already been deployed to detect harmful inputs in production at a small fraction of the cost of LLM-based classifiers \parencite{kramar2026probes}, and similar activation-based monitors could be combined with, or replace, entropy signals to detect communication and other unintended behaviours in a MAS more directly, although they would likewise require white-box access to the models. It is also necessary to study more broadly how tools, new information and other perturbations affect the entropy of a MAS, to establish non-communicating baselines across architectures and models, and to validate the approach in more realistic settings, with the aim of building a more robust detector of unintended behaviours and capabilities in general.
